\section{Task 3: Model Predictive Control (MPC)}

\subsection{Task Overview} In this notebook, we use \href{https://locuslab.github.io/mpc.pytorch/}{MPC PyTorch}, which is a fast and differentiable model predictive control solver for PyTorch. 
Our goal is to solve the \href{https://gymnasium.farama.org/environments/classic_control/pendulum/}{Pendulum} environment from \href{https://gymnasium.farama.org}{Gymnasium}, where we want to swing a pendulum to an upright position and keep it balanced there.

There are many helper functions and classes that provide the necessary tools for solving this environment using \textbf{MPC}. Some of these tools might be a little overwhelming, and that's fine, just try to understand the general ideas. Our primary objective is to learn more about \textbf{MPC}, not focusing on the physics of the pendulum environment.

On a final note, you might benefit from exploring the \href{https://github.com/locuslab/mpc.pytorch/}{source code} for \href{https://locuslab.github.io/mpc.pytorch/}{MPC PyTorch}, as this allows you to see how PyTorch is used in other contexts. To learn more about \textbf{MPC} and \textbf{mpc.pytorch}, you can check out \href{https://arxiv.org/abs/1703.00443}{OptNet} and \href{https://arxiv.org/abs/1810.13400}{Differentiable MPC}.

\textbf{Sections to be Implemented and Completed}

This notebook contains several placeholders (\texttt{TODO}) for missing implementations. In the final section, you can answer the questions asked in a markdown cell, which are the same as the questions in section \ref{sec:mpc-questions}.

\subsection{Questions}\label{sec:mpc-questions} 
You can answer the following questions in the notebook as well, but double-check to make sure you don't miss anything.
\subsubsection{Analyze the Results} 
Answer the following questions after running your experiments: 
\begin{itemize} 
    \item How does the number of LQR iterations affect the MPC? 

    \textbf{Answer:}

    The number of LQR (Linear Quadratic Regulator) iterations controls how thoroughly each MPC optimization step refines the action sequence. More iterations allow for better convergence toward an optimal control trajectory, improving performance and stability of the controller. However, increasing this number also increases computation time. If too few iterations are used, the control sequence may remain suboptimal, resulting in unstable or inefficient behavior.

    \item What if we didn't have access to the model dynamics? Could we still use MPC? 

    \textbf{Answer:}

    Without access to the true model dynamics, traditional MPC would not be directly applicable because it relies on simulating future states. However, model-free alternatives (like policy gradient methods) or learned dynamics models (using neural networks in model-based RL) can be used instead. In such cases, we can still use MPC by training a predictive model of the environment from data—this is often referred to as \textit{learned MPC} or \textit{model-based planning}.
    
    \item Do \texttt{TIMESTEPS} or \texttt{N\_BATCH} matter here? Explain. 

    \textbf{Answer:}

    Yes, both parameters are important:
    \begin{itemize}
        \item \texttt{TIMESTEPS} defines how many time steps into the future the MPC controller plans. A longer horizon allows for better long-term planning but increases computational load and may introduce more model error.
        \item \texttt{N\_BATCH} determines how many control sequences are optimized in parallel (i.e., batch size). A larger batch allows for better exploration of the control space and faster convergence of the optimal trajectory but also increases memory and compute requirements.
    \end{itemize}
    Both must be tuned based on the complexity of the environment and available compute resources.

    \item Why do you think we chose to set the initial state of the environment to the downward position? 
    
    Starting from the downward position (i.e., the pendulum hanging straight down) is the most challenging configuration for the task. The agent must learn a non-trivial control policy to swing the pendulum up and stabilize it—a classic benchmark problem for control. It ensures that the MPC controller is not merely maintaining balance, but actively planning a complex maneuver followed by stabilization.
    
    \item As time progresses (later iterations), what happens to the actions and rewards? Why

    As time progresses, the actions tend to become more subtle and smooth, focusing on maintaining the pendulum near the upright position. Initially, large actions are required to swing the pendulum up, resulting in lower rewards due to high energy usage and deviation from the goal. Once upright, the controller refines its actions to minimize energy and stabilize the pendulum, leading to increased and more consistent rewards over time. This transition reflects the success of MPC in both long-term planning and short-term precision control.

\end{itemize}